\begin{theorem}[Ramified Hasse preparation]
\label{mod:parameters-ramifiedhassepreparation}
Let $P$ be the prepared ramified prime-cyclic extension, let $H$ be its
stable-high odd parameter datum at depths $d,d'$, and fix an actual
downstairs stationary quotient representative $R$.  Put
\[
 \delta_F=\varpi_F^d,\qquad
 \delta_K=\varpi_K^{d'},\qquad
 \varpi_F=N_{K/F}(\varpi_K).
\]
The constructor \texttt{ramifiedHassePreparation} produces the complete
\texttt{RamifiedHasseComparisonData} for the actual lower phase
\texttt{H.downstairsOddStationaryPhase R} and actual upper phase
\texttt{H.upstairsOddStationaryPhase R}.  The upper denominator is the
ordered image of the lower denominator and its representative is literally
the image of $R$; neither is replaced by an abstract finite function.

For an integral lift $\tilde x$ set
\[
 Y_x=\varpi_K^{d'}\tilde x,\qquad
 T(x)=\overline{\operatorname{Tr}(Y_x)/\varpi_F^d},\qquad
 S(x)=\overline{E_2(Y_x)/\varpi_F^{2d}}.
\]
The local-geometry norm truncation, together with stationary
linearization for the exact quotient representative, proves the pointwise
transport formula
\[
 \phi_K(x)=\phi_F(T(x))\,\psi_0(-S(x)).
\]
Here $\phi_F,\phi_K$ are the genuine Lamprecht Hasse functions and
$\psi_0$ is the residual additive character derived from the same lower
representative and denominator.  Thus the negative second-symmetric
translation is retained before any finite normal form is chosen.

If the lower break is zero, the exact tame coordinate
$c=\overline{\varpi_K^{d'}/\varpi_F^d}$ satisfies
\[
 T(x)=\ell cx,\qquad
 S(x)=\binom{\ell}{2}c^2x^2,\qquad
 \bar\epsilon=\ell c^2.
\]
The first two identities come from reduction of the actual trace and
second elementary symmetric coefficient of the conjugates.  Hence
$\phi_K(c^{-1}x)=\phi_F(x)^\ell$.  In odd residue characteristic,
the lower Hasse function has the coefficient-one normal form
\[
 \phi_F(x)=\psi_0\left(\frac12x^2+\alpha x\right),
\]
and the upper function has the corresponding exact tame normal form with
quadratic coefficient $\ell/2$ and affine coefficient $\ell\alpha$.
The required quadratic-reciprocity sign is deduced from
$\ell\mid(|k_F|-1)$.  In residue characteristic two, the same actual
power relation is normalized by an explicit nonzero coordinate scale to a
\texttt{CharTwoRefinement}; the certificate records $|k_F|=2^f$, the
absolute-trace value at $1$, the difficult modulo-eight parity, and the
literal binomial correction.

If the break is positive, LocalGeometry makes $T(x)=0$.  Newton's
degree-two identity and its exact bilinear trace specialization give
\[
 -S(x)=\frac{\bar\epsilon}{2}x^2.
\]
The wild finite datum therefore has the exact upper normal form with zero
affine term.  Its field \texttt{epsilon\_eq} is filled directly by
\texttt{RamifiedHasseLocalGeometryData.wild\_epsilonResidue\_eq}:
\[
 \bar\epsilon=-\lambda^{\ell-1}.
\]
No critical-norm polynomial calculation or stable-to-critical trace
congruence is repeated in this node.

Finally the constructor copies the norm truncation, trace and
second-symmetric depths, epsilon normalization, tame and wild coordinate
origins, canonical residue equivalence, and trace specialization directly
from \texttt{RamifiedHasseLocalGeometryData}.  The resulting
\texttt{finiteComparison} is one of the three pointwise normal-form
certificates above, not an assumed phase comparison.
\LeanFile{LanglandsFirstMainLemma/Parameters/RamifiedHassePreparation.lean}
\lean{LanglandsFirstMainLemma.ramifiedHassePreparation}
\leanok
\SourceText{Theorem thm:ramified-Hasse-comparison.}
\uses{mod:parameters-ramifiedhasselocalgeometry,mod:parameters-ramifiedhassecomparison,mod:parameters-phasereduction,mod:parameters-high}
\FormalizationContract
The exported constructor takes the live \texttt{PrimeCyclicPreparation},
\texttt{HighStableOddParameterData}, and caller-supplied
\texttt{StationaryClassRepresentative}.  It fixes the source
uniformizer-power critical coordinates and returns the complete comparison
data for those literal phase wrappers.  The finite cases are proved
pointwise from the transported Lamprecht values; no final phase-power
identity is used as an input.
\end{theorem}
